\documentclass[11pt]{article}

\usepackage[a4paper,margin=1in]{geometry}
\usepackage{amsmath,amssymb,mathtools}
\usepackage{bm}
\usepackage{booktabs}
\usepackage{hyperref}
\usepackage{siunitx}
\usepackage{enumitem}

\title{Finite-Cell Gate Derivations and Remaining Operator Targets:\\
Surface Spectrum, Laplace Matching, GP/NLS Shell Response, and Hodge Capacity}

\author{Omar Iskandarani\\
Independent Researcher, Appingedam, The Netherlands}

\date{\today}

\newcommand{\LK}{\mathcal L_K}
\newcommand{\Harm}{\mathcal H}
\newcommand{\Kcell}{\mathcal K_{\rm cell}}
\newcommand{\Eeff}{E_{\rm eff}}
\newcommand{\EpNLS}{E_p^{\rm NLS}}
\newcommand{\alphacell}{\alpha_{\rm cell}}

\begin{document}
\maketitle

\begin{abstract}
The companion papers isolate a dimensionless fine-structure-scale closure based
on the ideal-trefoil ropelength, a spherical pressure cell, and a finite-cell
BEM/NLS stationary point.  This note records the updated status of the four
coefficient gates after the post-review calculations.  The pressure-sector
count is derived from the fixed-volume surface spectrum of a spherical cell:
\(k_\ell=\ell(\ell+1)-2\), so \(\ell=0\) is the compression sector,
\(\ell=1\) is the translation/centre-displacement sector, and \(\ell\ge2\) are
positive-stiffness shape modes; hence \(N_p=1+3=4\).  The radius action is
derived within the Laplace-matched reciprocal pressure model because the SST
kinetic and circular Laplace pressures agree,
\[
  \frac12\rho_{\rm core}\lVert\mathbf v_{\!\boldsymbol{\circlearrowleft}}\rVert^2
  =
  \frac{F_{\rm swirl}^{\max}}{2\pi r_c^2},
\]
giving \(A_\chi=\chi_R+N_p/\chi_R\) and \(\chi_R=2\).  The GP/NLS shell audit
derives \(\sigma=3+(2/3)w_\perp\); the value \(11/48\) follows for the
unweighted isotropic shell normalization \(w_\perp=1\), while deriving
\(w_\perp=1\) from the full microscopic core reduction remains open.  Finally,
the exterior Hodge phase calculation derives \(q_\phi=1\) and
\(\Lambda_\phi=4\pi R\mathcal K_{\rm cell}\), but the interior operator
identity \(\Lambda_\phi=E_{\rm eff}/2\) remains the required step for an
unconditional far-field stiffness.  Thus the present status is stronger than a
free closure model but still below a complete first-principles derivation of
the electromagnetic fine-structure constant.
\end{abstract}

\section{Purpose and status}

The preceding companion papers separate three levels of claim.  The first is
an algebraic closure audit.  The second is a spherical/NLS pressure-cell
construction.  The third is a finite-cell BEM/NLS stationary-point calculation
and conditional far-field consistency test.  The present note supplies the
missing logical bridge between a free closure model and a derived effective
model: it proves uniqueness of the minimal finite-cell variational reduction
under explicitly stated axioms.

The status of the result is therefore
\[
  \boxed{\text{derived within the unique minimal finite-cell variational reduction}.}
\]
It is not, by itself,
\[
  \boxed{\text{derived from the full microscopic dynamics without reduction assumptions}.}
\]
That stronger statement would require deriving the axioms below directly from
the underlying field theory.


\section{Post-review gate status}

The achieved results are:
\[
  N_p=4
  \quad [\text{derived in spherical surface spectrum; stable under perturbative, nonlinear star-shaped, and finite-perimeter global-topology checks}],
\]
\[
  A_\chi=\chi_R+\frac{N_p}{\chi_R},\quad \chi_R=2
  \quad [\text{derived in Laplace-matched reciprocal pressure model}],
\]
\[
  \frac{11}{48}
  \quad [\text{derived within unweighted isotropic GP/NLS shell reduction}],
\]
and
\[
  q_\phi=1,\quad \Lambda_\phi=4\pi R\mathcal K_{\rm cell}
  \quad [\text{derived by exterior Hodge capacity}].
\]
The remaining open targets are:
\[
  w_\perp=1
  \quad [\text{full microscopic GP/SST core reduction}],
\]
and
\[
  \Lambda_\phi=\frac{E_{\rm eff}}{2}
  \quad [\text{interior one-cell phase Hessian}].
\]
The earlier delay-stability route is no longer used as the primary justification
for \(N_p=4\), because its default parameters did not cleanly retain the
\(\ell=1\) sector.  The surface-spectrum route is parameter-free at quadratic
order and is therefore the preferred derivation.

\section{Axioms of the minimal finite-cell reduction}

Let the finite cell be represented by a spherical boundary \(S^2\), with
outward unit normal \(\mathbf n\), pressure field \(p\), dimensionless shell
thickness \(\eta=a/R_{\rm cell}\), cell-radius ratio
\(\chi_R=R_{\rm cell}/L_K\), and exterior phase field \(u\).  The minimal
quadratic reduction is assumed to satisfy:
\begin{enumerate}[label=A\arabic*.]
\item \(SO(3)\) covariance on the spherical boundary.
\item Locality at quadratic order on \(S^2\).
\item Positive quadratic energy for stable boundary perturbations.
\item Pressure constraints only for total compression and centre displacement:
\[
  \int_{S^2}p\,d\Omega=P_0,\qquad
  \int_{S^2}p\,\mathbf n\,d\Omega=\mathbf P_1.
\]
\item Isotropic NLS/GP shell-gradient response at second order.
\item Reciprocal inner/outer pressure balance at leading scale order.
\item Leading far-field coupling carried only by the exterior \(1/r\) harmonic
phase mode.
\end{enumerate}

These assumptions define the class in which the uniqueness statement is made.

\section{The finite-cell effective action}

At quadratic order the effective action separates into
\[
  \mathcal A_{\rm eff}
  =
  \mathcal A_p[p]
  +
  \mathcal A_{\rm shell}[\eta]
  +
  A_\chi(\chi_R)
  +
  \mathcal A_\phi[u].
\]
The four pieces are fixed below by the axioms, up to an overall normalization
that cancels in the dimensionless coefficient gates.

\section{Theorem 1: constrained pressure minimization gives \(N_p=4\)}

Consider
\[
  \mathcal A_p[p]
  =
  \frac12
  \int_{S^2}
  |\nabla_{S^2}p|^2\,d\Omega
\]
subject to
\[
  \int_{S^2}p\,d\Omega=P_0,
  \qquad
  \int_{S^2}p\,\mathbf n\,d\Omega=\mathbf P_1.
\]
Introducing Lagrange multipliers \(\lambda_0\) and
\(\boldsymbol{\lambda}_1\), the Euler--Lagrange equation is
\[
  -\Delta_{S^2}p
  =
  \lambda_0+\boldsymbol{\lambda}_1\cdot\mathbf n.
\]
The constant function belongs to \(\Harm_0\), and the coordinate functions
\(n_x,n_y,n_z\) span \(\Harm_1\).  Hence the minimal pressure subspace is
\[
  \mathcal P_{\rm min}
  =
  \Harm_0\oplus\Harm_1.
\]
Since
\[
  \dim\Harm_\ell=2\ell+1,
\]
one obtains
\[
  N_p
  =
  \dim\Harm_0+\dim\Harm_1
  =
  1+3
  =
  4.
\]
This proves uniqueness of \(N_p=4\) inside the lowest \(SO(3)\)-closed pressure
subspace satisfying the volume and centre-displacement constraints.  Modes
\(\ell\ge2\) represent quadrupolar and higher shape deformations and are not
required by the constraints.


\section{Theorem 1b: fixed-volume surface spectrum gives the same \(N_p=4\)}

The constrained pressure argument above is strengthened by the fixed-volume
surface spectrum.  Let the spherical boundary be
\[
  r(\Omega)=R[1+\epsilon f(\Omega)].
\]
At fixed volume,
\[
  \delta^2 A
  =
  \frac{R^2}{2}
  \int_{S^2}
  \left(
  |\nabla_{S^2}f|^2-2f^2
  \right)d\Omega.
\]
For \(f=Y_{\ell m}\),
\[
  k_\ell
  =
  \ell(\ell+1)-2
  =
  (\ell-1)(\ell+2).
\]
Thus \(k_1=0\), identifying the three translation/centre-displacement modes.
For \(\ell\ge2\), \(k_\ell>0\), so these are positive-stiffness shape modes
integrated out at zeroth order.  The \(\ell=0\) mode is retained separately as
the scalar compression degree of freedom.  Hence
\[
  \mathcal P_{\rm min}
  =
  \mathcal H_0\oplus\mathcal H_1,
  \qquad
  N_p=1+3=4.
\]
This is the preferred derivation of \(N_p=4\) in the present package.


\section{Theorem 1c: finite-core and nonlinear \(\ell\ge2\) shape stability}

The first positive fixed-volume shape-mode gap is
\[
  k_2=4.
\]
Let \(\delta H\) be the finite-core or weak nonspherical perturbation of the
\(\ell\ge2\) shape Hessian.  If
\[
  \|\delta H\|_{\rm op}<\frac{k_2}{2}=2,
\]
then the \(\ell\ge2\) block remains separated from the retained
\(\mathcal H_0\oplus\mathcal H_1\) pressure manifold.  The finite-core scale in
the trefoil cell is
\[
  \eta_K=\frac{1}{4\mathcal L_K}=1.52703116982\times10^{-2},
\]
far below this half-gap bound.

A nonlinear star-shaped audit tests finite-amplitude deformations
\[
  r(\Omega)=
  \exp\!\left(
  \sum_{\ell=2}^{\ell_{\max}}\sum_m a_{\ell m}Y_{\ell m}(\Omega)
  \right)
\]
with exact volume rescaling to \(V=4\pi/3\).  For \(\ell_{\max}=6\), the
\(\ell\ge2\) coefficient space has 45 real modes.  Across 300 random samples
per amplitude at \(\eta_K,0.05,0.10,0.25,0.50,1.00\), no negative fixed-volume
area excess was found.  At the physical scale,
\[
  \min\Delta A_{\rm random}=1.71825494947\times10^{-3}.
\]
Moreover, 20 L-BFGS-B local minimizations from amplitude-one random starts
returned to the spherical solution, with
\[
  \Delta A_{\rm final}\lesssim 1.2\times10^{-12},
  \qquad
  \max\lVert a_{\ell m}^{\rm final}\rVert
  =
  5.830157\times10^{-7}.
\]
Thus nonlinear \(\ell\ge2\) star-shaped deformations do not enter the leading
pressure manifold in the tested class.

\section{Theorem 1d: global topology admissibility}

The nonlinear audit covers star-shaped surfaces.  The global loophole is closed
for the larger admissible class of bounded finite-perimeter pressure cells by
the Euclidean isoperimetric inequality:
\[
  P(\Omega)^3\ge 36\pi V(\Omega)^2,
\]
where \(P(\Omega)\) is the perimeter and \(V(\Omega)\) the enclosed volume.
Equality holds only for a ball, up to null sets.  Hence, at fixed
\(V(\Omega)=4\pi R^3/3\),
\[
  P(\Omega)\ge 4\pi R^2,
  \qquad
  \Delta A=P(\Omega)-4\pi R^2\ge0.
\]
Therefore no admissible non-star-shaped, disconnected, or higher-genus
finite-perimeter pressure cell can lower the fixed-volume area below the
sphere.  Self-intersecting immersed surfaces are not admissible reduced
boundaries of a physical pressure cell; if counted with multiplicity, their
area cannot improve the isoperimetric bound.  Thus global topology does not
promote \(\ell\ge2\) modes into the leading pressure manifold.

\section{Theorem 2: isotropic NLS/GP shell variation gives \(\sigma=11/3\)}

Let
\[
  \eta=\frac{a}{R_{\rm cell}}.
\]
The volume Jacobian for a relative inward shell displacement is
\[
  \frac{V(R-a)}{V(R)}
  =
  (1-\eta)^3
  =
  1-3\eta+3\eta^2+O(\eta^3).
\]
After the first-order displacement is absorbed into the equilibrium radius, the
second-order volume contribution is
\[
  \sigma_{\rm vol}=3.
\]
The isotropic NLS/GP gradient term contributes through the transverse
projector
\[
  P_\perp=I-\hat{\mathbf r}\hat{\mathbf r}.
\]
Using
\[
  \langle \hat r_i\hat r_j\rangle_{S^2}
  =
  \frac13\delta_{ij},
\]
one gets
\[
  \langle P_\perp\rangle_{S^2}
  =
  \frac23 I.
\]
Equivalently,
\[
  \langle\sin^2\theta\rangle_{S^2}
  =
  \frac{\int_0^\pi \sin^2\theta\,\sin\theta\,d\theta}
       {\int_0^\pi \sin\theta\,d\theta}
  =
  \frac23.
\]
Thus the second-order shell coefficient is
\[
  \sigma
  =
  \sigma_{\rm vol}+\sigma_\perp
  =
  3+\frac23
  =
  \frac{11}{3}.
\]


The GP core-profile audit confirms the angular result and computes the radial
profile.  It shows that the normalized angular factor \(2/3\) is profile
independent for isotropic radial weights.  It also shows that alternative
profile-weight proxies do not generically reproduce \(11/48\).  Therefore the
precise status is
\[
  \sigma=3+\frac23w_\perp,
\]
with \(11/48\) obtained for \(w_\perp=1\).  The equality \(w_\perp=1\) is the
remaining microscopic normalization gate.

\section{Theorem 3: reciprocal pressure stationarity gives \(\chi_R=2\)}


The equal inner/outer normalization is now supported by the SST pressure
identity.  With the canonical constants,
\[
  P_{\rm kin}
  =
  \frac12\rho_{\rm core}
  \lVert\mathbf v_{\!\boldsymbol{\circlearrowleft}}\rVert^2
  =
  2.329244600448\times10^{30}\ {\rm Pa},
\]
and
\[
  P_{\rm Laplace}
  =
  \frac{F_{\rm swirl}^{\max}}{2\pi r_c^2}
  =
  2.329244600448\times10^{30}\ {\rm Pa}.
\]
Therefore the leading pressure ratio is
\[
  \mu=\frac{P_{\rm in}}{P_{\rm out}}=1.
\]
The general reciprocal action
\[
  A_\chi=\chi_R+\mu\frac{N_p}{\chi_R}
\]
then reduces to the self-dual form used below.

Let
\[
  \chi_R=\frac{R_{\rm cell}}{L_K}.
\]
The leading outer pressure-extension cost is proportional to \(\chi_R\), while
the inner crowding/compliance cost is reciprocal and proportional to the number
of pressure sectors.  Inner/outer reciprocity therefore fixes the minimal
scale action to
\[
  A_\chi(\chi_R)
  =
  \chi_R+\frac{N_p}{\chi_R}.
\]
Stationarity gives
\[
  \frac{dA_\chi}{d\chi_R}
  =
  1-\frac{N_p}{\chi_R^2}
  =
  0,
\]
so
\[
  \chi_R=\sqrt{N_p}=2.
\]
Equivalently, if
\[
  A_\chi=\chi_R+\frac{\lambda_\chi}{\chi_R},
\]
then
\[
  \lambda_\chi=N_p=4.
\]
This is unique in the leading-order class of reciprocal, scale-homogeneous
inner/outer pressure actions.  If additional terms such as \(\chi_R^2\),
\(\chi_R^{-2}\), or \(\log\chi_R\) are admitted at the same order, uniqueness is
lost; such terms are excluded by the minimal scale-reciprocity axiom.

\section{Theorem 4: the exterior phase Hessian fixes \(q_\phi=1\)}

An exterior harmonic field has the multipole expansion
\[
  u(r,\theta,\phi)
  =
  \sum_{\ell,m}
  a_{\ell m}
  \frac{Y_{\ell m}(\theta,\phi)}{r^{\ell+1}}.
\]
Only \(\ell=0\) gives a \(1/r\) field.  The \(\ell=1\) sector begins at
\(r^{-2}\), and all \(\ell\ge2\) modes are shorter-ranged.  Therefore the
leading far-field charge is the constant boundary phase, i.e. the unique
\(H^0(S^2)\) mode.  Since \(S^2\) is connected,
\[
  q_\phi=\dim H^0(S^2)=1.
\]
For the exterior monopole
\[
  u(r)=\frac{\phi}{r},
\]
one has
\[
  \int_{r\ge1}
  |\nabla u|^2\,d^3x
  =
  \int_1^\infty
  4\pi r^2\frac{\phi^2}{r^4}\,dr
  =
  4\pi\phi^2.
\]
Thus
\[
  \Lambda_\phi=4\pi\Kcell.
\]
The unit \(H^0\) one-cell phase Hessian gives
\[
  \Lambda_\phi=\frac{\Eeff}{2},
\]
and therefore
\[
  \Kcell
  =
  \frac{\Eeff}{8\pi}.
\]


The exterior Hodge calculation derives \(q_\phi=1\) and the capacity relation
\(\Lambda_\phi=4\pi R\mathcal K_{\rm cell}\).  It does not by itself derive the
interior identity \(\Lambda_\phi=E_{\rm eff}/2\).  Thus
\[
  \mathcal K_{\rm cell}=\frac{E_{\rm eff}}{8\pi R}
\]
is conditional on the interior one-cell phase-Hessian result.  For \(R=1\), the
conditional target is \(\mathcal K_{\rm cell}=E_{\rm eff}/(8\pi)\).

\section{Resulting pressure scale}

Combining Theorems 1--3,
\[
  E_p^{(0)}
  =
  N_p\frac{4\pi}{3}\LK
  =
  \frac{16\pi}{3}\LK.
\]
Moreover,
\[
  \eta_K
  =
  \frac{1}{2\chi_R\LK}
  =
  \frac{1}{4\LK}.
\]
With \(\sigma=11/3\),
\[
  \sigma\eta_K^2
  =
  \frac{11}{3}\frac{1}{16\LK^2}
  =
  \frac{11}{48\LK^2}.
\]
Therefore
\[
  \boxed{
  \EpNLS
  =
  \frac{16\pi}{3}\LK
  \left(
  1-\frac{11}{48\LK^2}
  \right).
  }
\]
Theorem 4 then gives the far-field stiffness relation
\[
  \Kcell=\frac{\Eeff}{8\pi},
\]
so that the leading two-cell interaction has coefficient
\[
  \alphacell
  =
  \frac{1}{4\pi\Kcell}
  =
  \frac{2}{\Eeff}.
\]


\section{Reproducibility map}

The supporting scripts and output files are stored in \texttt{../code/} relative
to the \texttt{Manuscripts} directory:
\begin{itemize}
\item \texttt{derive\_pressure\_mode\_cutoff\_surface\_spectrum.py}:
linear surface-spectrum derivation of \(N_p=4\);
\item \texttt{test\_finite\_core\_nonspherical\_shape\_corrections.py}:
perturbative finite-core/nonspherical half-gap audit;
\item \texttt{solve\_nonlinear\_shape\_stability.py}:
nonlinear star-shaped \(\ell\ge2\) shape-stability solve;
\item \texttt{derive\_gp\_core\_profile\_second\_variation.py}:
GP/NLS shell coefficient and \(w_\perp\) sensitivity audit;
\item \texttt{derive\_pressure\_self\_duality\_from\_laplace\_matching.py}:
Laplace-matched reciprocal radius action;
\item \texttt{solve\_one\_cell\_hodge\_phase\_hessian.py}:
exterior Hodge capacity and \(q_\phi=1\).
\end{itemize}
The numerical CSV and Markdown outputs are stored in the corresponding
\texttt{../code/outputs\_*} directories and summarized in
\texttt{../code/DERIVED\_GATE\_EVIDENCE\_MANIFEST.json}.

\section{Falsification and robustness checks}

The uniqueness theorem depends on the minimal axioms.  The following checks
test those axioms rather than the algebra:
\begin{enumerate}
\item include the \(\ell=2\) pressure sector and verify that it enters only as a
higher-order shape correction;
\item replace \(P_\perp\) by a weighted transverse projector and verify that
isotropic NLS/GP reduction fixes the weight to unity;
\item replace \(A_\chi=\chi_R+N_p/\chi_R\) by
\(a\chi_R+bN_p/\chi_R\) and verify that inner/outer reciprocity enforces
\(a=b\);
\item compute the one-cell phase Hessian directly from the full finite-cell
operator and verify \(\Lambda_\phi=\Eeff/2\).
\end{enumerate}

\section{Conclusion}

The four gate calculations now have differentiated status.  The pressure-sector
count is derived by the spherical fixed-volume surface spectrum and is stable
under perturbative finite-core, nonlinear star-shaped, and global
finite-perimeter topology checks:
\[
  k_\ell=(\ell-1)(\ell+2),
  \qquad
  \mathcal P_{\rm min}=\mathcal H_0\oplus\mathcal H_1,
  \qquad
  N_p=4.
\]
The radius action is derived within the Laplace-matched reciprocal pressure
model because the SST kinetic and circular Laplace pressures agree, giving
\(\mu=1\), \(A_\chi=\chi_R+N_p/\chi_R\), and \(\chi_R=2\).  The finite-shell
coefficient \(11/48\) is derived within the unweighted isotropic GP/NLS shell
reduction, but the primitive-equation derivation of \(w_\perp=1\) remains open.
The exterior phase-Hodge calculation derives \(q_\phi=1\) and
\(\Lambda_\phi=4\pi R\mathcal K_{\rm cell}\), while the interior Hessian
identity \(\Lambda_\phi=E_{\rm eff}/2\) remains open.

Thus the package has moved beyond a free closure model for \(N_p\) and
\(\chi_R\), and it has isolated \(11/48\) to a single microscopic normalization
gate.  A complete first-principles derivation of the electromagnetic
fine-structure constant still requires the two remaining primitive-equation
results:
\[
  w_\perp=1,
  \qquad
  \Lambda_\phi=\frac{E_{\rm eff}}{2}.
\]

\begin{thebibliography}{9}

\bibitem{MorseFeshbach1953}
P. M. Morse and H. Feshbach,
\emph{Methods of Theoretical Physics},
McGraw--Hill, New York (1953).

\bibitem{ArfkenWeberHarris2013}
G. B. Arfken, H. J. Weber, and F. E. Harris,
\emph{Mathematical Methods for Physicists}, 7th ed.,
Academic Press, Oxford (2013).
ISBN: 978-0123846549.

\bibitem{Pitaevskii1961}
L. P. Pitaevskii,
``Vortex Lines in an Imperfect Bose Gas,''
\emph{Soviet Physics JETP} \textbf{13}, 451--454 (1961).

\bibitem{RobertsGrant1971}
P. H. Roberts and J. Grant,
``Motions in a Bose Condensate. I. The Structure of the Large Circular Vortex,''
\emph{Journal of Physics A: General Physics} \textbf{4}, 55--72 (1971).
DOI: 10.1088/0305-4470/4/1/306.

\bibitem{Saffman1992}
P. G. Saffman,
\emph{Vortex Dynamics},
Cambridge University Press, Cambridge (1992).
ISBN: 978-0521477390.

\bibitem{Evans2010}
L. C. Evans,
\emph{Partial Differential Equations}, 2nd ed.,
American Mathematical Society, Providence, RI (2010).
ISBN: 978-0821849743.

\bibitem{Osserman1978}
R. Osserman,
``The isoperimetric inequality,''
\emph{Bulletin of the American Mathematical Society} \textbf{84},
1182--1238 (1978).
DOI: 10.1090/S0002-9904-1978-14553-4.

\bibitem{EvansGariepy2015}
L. C. Evans and R. F. Gariepy,
\emph{Measure Theory and Fine Properties of Functions},
revised ed.,
CRC Press, Boca Raton (2015).
ISBN: 978-1482242386.

\bibitem{Maggi2012}
F. Maggi,
\emph{Sets of Finite Perimeter and Geometric Variational Problems},
Cambridge University Press, Cambridge (2012).
DOI: 10.1017/CBO9781139108133.

\end{thebibliography}

\end{document}
